% Generated by scripts/build_reports.py; edit reviewed Markdown inputs.
% Compile with XeLaTeX.
% Options for packages loaded elsewhere
\PassOptionsToPackage{unicode}{hyperref}
\PassOptionsToPackage{hyphens}{url}
\PassOptionsToPackage{dvipsnames,svgnames,x11names}{xcolor}
\documentclass[
  11pt,
]{article}
\usepackage{xcolor}
\usepackage[margin=1in]{geometry}
\usepackage{amsmath,amssymb}
\setcounter{secnumdepth}{-\maxdimen} % remove section numbering
\usepackage{iftex}
\ifPDFTeX
  \usepackage[T1]{fontenc}
  \usepackage[utf8]{inputenc}
  \usepackage{textcomp} % provide euro and other symbols
\else % if luatex or xetex
  \usepackage{unicode-math} % this also loads fontspec
  \defaultfontfeatures{Scale=MatchLowercase}
  \defaultfontfeatures[\rmfamily]{Ligatures=TeX,Scale=1}
\fi
\usepackage{lmodern}
\ifPDFTeX\else
  % xetex/luatex font selection
\fi
% Use upquote if available, for straight quotes in verbatim environments
\IfFileExists{upquote.sty}{\usepackage{upquote}}{}
\IfFileExists{microtype.sty}{% use microtype if available
  \usepackage[]{microtype}
  \UseMicrotypeSet[protrusion]{basicmath} % disable protrusion for tt fonts
}{}
\makeatletter
\@ifundefined{KOMAClassName}{% if non-KOMA class
  \IfFileExists{parskip.sty}{%
    \usepackage{parskip}
  }{% else
    \setlength{\parindent}{0pt}
    \setlength{\parskip}{6pt plus 2pt minus 1pt}}
}{% if KOMA class
  \KOMAoptions{parskip=half}}
\makeatother
\usepackage{longtable,booktabs,array}
\usepackage{caption}
\captionsetup[table]{skip=6pt}
\newcounter{none} % for unnumbered tables
\usepackage{calc} % for calculating minipage widths
% Correct order of tables after \paragraph or \subparagraph
\usepackage{etoolbox}
\makeatletter
\patchcmd\longtable{\par}{\if@noskipsec\mbox{}\fi\par}{}{}
\makeatother
% Allow footnotes in longtable head/foot
\IfFileExists{footnotehyper.sty}{\usepackage{footnotehyper}}{\usepackage{footnote}}
\makesavenoteenv{longtable}
\setlength{\emergencystretch}{3em} % prevent overfull lines
\providecommand{\tightlist}{%
  \setlength{\itemsep}{0pt}\setlength{\parskip}{0pt}}
\usepackage{amsmath,amssymb}
\usepackage{xurl}
\usepackage{fancyhdr}
\pagestyle{fancy}
\fancyhf{}
\fancyhead[L]{\small Pure Thought | Revised research specifications}
\fancyfoot[L]{\small 2026-09-17}
\fancyfoot[R]{\small\thepage}
\setlength{\headheight}{15pt}
\setlength{\emergencystretch}{3em}
\allowdisplaybreaks
\usepackage{bookmark}
\IfFileExists{xurl.sty}{\usepackage{xurl}}{} % add URL line breaks if available
\urlstyle{same}
% fallback for those not using the hyperref driver hyperxmp:
\makeatletter
\@ifundefined{xmpquote}{\newcommand{\xmpquote}[1]{#1}}{}
\makeatother
\hypersetup{
  pdftitle={S10. Certified periodic orbits of a specified three-body model},
  pdfauthor={Pure Thought research portfolio},
  colorlinks=true,
  linkcolor={blue},
  filecolor={Maroon},
  citecolor={Blue},
  urlcolor={blue},
  pdfcreator={LaTeX via pandoc}}

\title{S10. Certified periodic orbits of a specified three-body model}
\author{Pure Thought research portfolio}
\date{Revised 17 September 2026}

\begin{document}
\maketitle
\thispagestyle{empty}
\clearpage

\section{S10. Certified periodic orbits of a specified three-body model}

\textbf{Source:} synthesis challenge 27;
\href{https://github.com/rick-stevens-ai/Pure-Thought/blob/main/NOTES/30-Most-Compelling-Pure-Thought-AI-Challenges.md}{original
synthesis}. This topic is distinct from, or adds a distinct objective
to, the canonical PRD catalog.
\href{https://github.com/rick-stevens-ai/Pure-Thought/blob/main/reviews/SYNTHESIS-CROSSWALK.md}{Crosswalk}

\subsection{Plain-language guide}

\subsubsection{The problem in everyday terms}

A simulated three-body trajectory may nearly return to where it started,
but a small numerical error can disguise whether an exactly repeating
orbit exists. This problem converts a promising numerical trajectory
into a proof that a nearby periodic solution exists and stays away from
collisions, then examines how small disturbances behave over one cycle.

\subsubsection{Key terms}

\begin{itemize}
\item
  \textbf{Three-body problem} --- The motion of three mutually
  gravitating bodies in the specified Newtonian model.
\item
  \textbf{Restricted model} --- A version in which one body has
  negligible mass and does not affect the other two.
\item
  \textbf{Periodic orbit} --- A trajectory that returns to the same
  state after a fixed positive time.
\item
  \textbf{Shooting method} --- Adjusting an initial state until its
  evolved state satisfies the return conditions.
\item
  \textbf{Phase condition} --- A rule fixing the arbitrary starting
  point along a periodic trajectory.
\item
  \textbf{Validated integration} --- Numerical evolution with guaranteed
  enclosures for errors.
\item
  \textbf{Collision exclusion} --- A proved positive lower bound on
  relevant body separations along the orbit.
\item
  \textbf{Floquet multiplier} --- An eigenvalue describing linear
  perturbation growth or rotation after one orbit.
\item
  \textbf{Neutral direction} --- A direction with constrained or trivial
  linear behavior arising from symmetries or conservation laws.
\end{itemize}

\subsubsection{Why this matters}

Periodic orbits are organizing reference motions in gravitational
dynamics. Certifying their existence gives subsequent studies of nearby
motion a dependable starting point and separates exact solutions from
visual patterns in a simulation.

\subsubsection{What progress would mean}

Progress would enclose an orbit and its period, exclude collisions, and
classify the reduced linear stability when the bounds permit. Linear
stability would not alone prove nonlinear or indefinite stability, and a
periodic-orbit proof would not establish connecting trajectories between
different orbits.

\subsection{Revised core question}

Can a numerically proposed periodic orbit be enclosed by a validated
shooting proof and classified by its reduced Floquet spectrum?

\subsection{Precise scope, target and outputs}

Choose one nondimensional Newtonian three-body model with positive
rational masses, or a planar circular restricted model with fixed mass
ratio. Declare collision exclusions and conserved quantities. Define a
Poincaré section or phase condition to remove time translation, and
reduce continuous symmetries. Solve the resulting shooting residual
using a validated flow integrator and interval Newton/Krawczyk bounds
for initial conditions and period. Integrate variational equations with
enclosures for the monodromy matrix; remove the expected neutral
directions before interpreting multipliers.

\subsection{Detailed critique of the synthesis framing}

The original bundles orbit existence, complete family discovery, linear
stability, invariant manifolds and heteroclinic connections into one
certificate. A small shooting residual is not existence, and
stable-looking Floquet multipliers are not nonlinear stability.
Near-return trajectories do not establish a connection between invariant
manifolds. The revised first target is a single orbit with a phase-fixed
existence proof; continuation, manifold enclosure and connection proofs
are explicit later stages. This also preserves the figure-eight idea
incorrectly folded into PRD 27's central configurations.

\subsection{Validation and milestones}

Baseline: certify one supplied orbit and its period. Strong: validated
continuation over a parameter interval and a reduced Floquet
classification. A manifold connection requires additional
covering/intersection proof, not a plot of approaching trajectories.

Use the shared
\href{https://github.com/rick-stevens-ai/Pure-Thought/blob/main/reviews/EVIDENCE-STANDARD.md}{evidence
standard}. Unsuccessful search is unresolved unless a complete exclusion
proof is supplied. These are proposed baseline specifications, not
completed scientific results.

\subsection{Displayed formulation and claim boundary}

\href{https://github.com/rick-stevens-ai/Pure-Thought/blob/main/supplemental/S10.md}{Supplemental
specification}.

\[F(z_0,T)=\varphi_T(z_0)-z_0=0,\qquad
M=D\varphi_T(z_0).\]

The shooting residual must first be augmented and reduced using a phase
condition, conserved quantities and continuous symmetries so the
validated root problem has the correct dimension. Interval flow and root
bounds establish existence; the reduced Floquet multipliers classify
linear behavior. Neither a near return nor stable-looking multipliers
prove nonlinear stability or a manifold connection.

\section{Evidence and execution standard}

This standard applies to the revised specifications. It is a proposed
contract for future research runs; this editorial revision does not
supply the scientific certificates themselves.

\subsection{What ``pure thought'' means here}

Use explicitly stated mathematical models and deductions without fitting
conclusions to experimental outcomes. Synthetic examples, exact model
inputs, analytic references, numerical search, and established software
are allowed. Physical constants, electronic-structure approximations,
constitutive laws, and empirical parameter sets must be identified when
used. A model-derived result is conditional on that model.

The original ban on legacy software conflicts with its proposed reliance
on NumPy, SciPy, CVXPY, nauty, and other libraries. Reimplementing every
optimizer and algebra system is not necessary for scientific
independence. Prefer a small independently checkable certificate and a
documented trusted base. A fresh implementation can be a project
objective, but must not be confused with a mathematical proof or novelty
claim.

\subsection{Four evidence levels}

{\def\LTcaptype{none} % do not increment counter
\begin{longtable}[]{@{}
  >{\raggedright\arraybackslash}p{(\linewidth - 4\tabcolsep) * \real{0.3333}}
  >{\raggedright\arraybackslash}p{(\linewidth - 4\tabcolsep) * \real{0.3333}}
  >{\raggedright\arraybackslash}p{(\linewidth - 4\tabcolsep) * \real{0.3333}}@{}}
\toprule\noalign{}
\begin{minipage}[b]{\linewidth}\raggedright
Label
\end{minipage} & \begin{minipage}[b]{\linewidth}\raggedright
What it establishes
\end{minipage} & \begin{minipage}[b]{\linewidth}\raggedright
What it does not establish
\end{minipage} \\
\midrule\noalign{}
\endhead
\bottomrule\noalign{}
\endlastfoot
Numerical evidence & A reproducible computation at stated precision,
sample size, and discretization & Exact truth, continuum coverage, or a
universal theorem \\
Finite-model certificate & An exact/validated statement about the
specified finite model or finite relaxation & An omitted tail,
thermodynamic limit, or physical realization \\
Conditional mathematical theorem & A proof over the stated domain with
every hypothesis and error controlled & Truth outside those assumptions
or a demonstrated physical system \\
Formal proof & A checked theorem in a named proof assistant, with
imported axioms/dependencies listed & Correct physical modeling merely
by virtue of formalization \\
\end{longtable}
}

Formalization is an optional extra deliverable unless the specific
project makes it mandatory. High precision, JSON output, solver status
and small residuals do not by themselves confer any of these proof
levels.

\subsection{Claim directions}

For minimization over a larger feasible set, the optimum is a lower
bound on the original minimum. A physical feasible point supplies an
upper bound. For maximization the directions reverse. Track feasibility
of the primal/dual witnesses separately from their objective values.

Infeasibility of an inner ansatz excludes only that ansatz.
Infeasibility of an outer relaxation can exclude the original problem if
the map into the relaxation is proved and all approximations preserve
necessity. Passing necessary tests means ``not excluded''; existence
needs a constructive sufficient argument.

A smallest observed example is not a minimality theorem. A complete
catalog must define the universe, equivalence relation, boundaries, and
proof that pruning misses nothing. A finite-size numerical trend does
not prove an asymptotic limit.

\subsection{Minimum run manifest}

Before calculation, record the problem ID and revision hash, exact
inputs and conventions, quantified domain, equivalence relation, target
observable, accepted evidence level, tolerances, and resource limits.
Record dependencies and precision/rounding mode. Where randomization is
used, preserve seeds, trial counts and raw sufficient statistics.
Resource limits are planning decisions, not evidence that a mathematical
solution cannot exist.

An output bundle should include:

\begin{itemize}
\tightlist
\item
  The statement being claimed, including all assumptions and
  quantifiers.
\item
  Construction or proof object, plus an independently runnable verifier
  and its trust assumptions.
\item
  Numerical/error ledger separating roundoff, discretization,
  truncation, optimization, sampling and modeling error where relevant.
\item
  Positive reference cases and negative controls that the verifier must
  reject.
\item
  Coverage: included domain, omitted domain, unresolved regions and
  termination reason.
\end{itemize}

These are requirements for future implementations, not a promise of
polynomial-size certificates for every problem.

\subsection{Checking specific certificate types}

\textbf{LP/SDP/SOS:} verify equality constraints, objective direction,
and cone feasibility using exact arithmetic or outward-rounded
intervals. Rationalizing a numerical solution can break PSD constraints.
Correct residuals with justified norm bounds; do not silently project
and reuse the old objective. Distinguish a proof of a polynomial
approximation from a proof for the original function.

\textbf{Interval numerics:} state function enclosures, domains and
outward rounding. A small root box needs existence and, if claimed,
uniqueness. Cover the complement for a global enumeration. Include
singularities and tails rather than discarding them numerically.

\textbf{SAT/SMT:} a solver's UNSAT label is not a certificate. Preserve
a proof supported by a checker and the proof that the encoding, symmetry
reduction and blocked solutions match the mathematical universe. DRAT
applies to an appropriate Boolean CNF proof, not automatically to every
real-arithmetic or SDP claim.

\textbf{Statistics:} specify the estimand, dependence assumptions,
confidence procedure and bias. Zero observed failures does not prove
zero failure probability. MCMC error bars based on estimated
autocorrelation are not generally rigorous mixing guarantees.

\subsection{Milestones and novelty}

Baseline means a working, independently validated result on a concrete
reference problem. Strong means a broader or sharper result with its
claimed evidence level. Research extension means a possible novel result
whose feasibility and significance need review. A negative theorem can
be valuable; an unsuccessful search can also be useful if its scope is
documented, but the two are not interchangeable.

Do not promise ``publication-quality'' in a fixed number of months.
Before claiming novelty, conduct a focused current literature comparison
against the exact theorem/model/metric. The sources in this revision
check key definitions and limitations; they are not an exhaustive
novelty audit across all 30 fields.

\end{document}
